% !TEX root = z-main.tex

The project aimed to digitally preserve and disseminate Guatemalan cultural heritage through the development of three-dimensional models of the archaeological site of Kaminaljuyú. This site, of great historical significance for the central highlands of Guatemala, is largely buried beneath urban development. This situation has hindered its access and study.

To address this issue, six three-dimensional models of selected mounds were created based on archaeological information validated by specialists. The work made it possible to digitally reconstruct structures that are no longer visible today and to prepare these models for integration into an augmented reality application oriented toward education and cultural dissemination.

Faithful three-dimensional representations were obtained. These incorporated textures and finishes that enhanced the understanding of the original appearance of the mounds. The final files were stored in compatible formats, ensuring their preservation and potential reuse in future research or dissemination projects.

It was concluded that the digitalization of archaeological sites through three-dimensional modeling and its application in augmented reality environments constitutes an effective strategy for heritage conservation. Its continuation is suggested in other archaeological areas of the country.